\section{Détection et protection}\label{sec:detection-protection}

Suricata est exécuté sur la passerelle et observe les requêtes HTTP avant leur traduction d'adresse.
Il écrit ses alertes dans un journal partagé avec l'hôte, ce qui permet de relier un rejeu du scénario à une alerte horodatée.
Deux règles locales sont dédiées aux scénarios : l'une détecte au moins cinq requêtes vers le formulaire de force brute en soixante secondes ; l'autre recherche une apostrophe encodée dans la ressource SQLi.
Par exemple, la règle SQLi est la suivante :

\begin{quote}
    \small\ttfamily
    alert http any any -> 172.28.20.20 80 (msg:\textquotedbl{}LAB DVWA SQL injection pattern\textquotedbl;\\
    \quad flow:to\_server,established; http.uri.raw; content:\textquotedbl{}/vulnerabilities/sqli/\textquotedbl;\\
    \quad content:\textquotedbl{}\%27\textquotedbl; nocase; sid:1000002; rev:1;)
\end{quote}

Ces signatures sont volontairement spécifiques à la maquette.
Elles doivent être évaluées face aux faux positifs et aux faux négatifs, notamment pour les autres encodages d'une charge SQLi.
Le choix est justifié à partir des journaux collectés et de la documentation Suricata~\cite{suricata}.
La protection de base de la passerelle applique une politique restrictive : les flux transférés sont refusés par défaut, puis seul le trafic HTTP de Kali vers DVWA est autorisé.
Après une alerte validée, le script de démonstration insère une règle iptables isolant l'adresse Kali.
Cette réponse est délibérément manuelle, explicite et réversible ; elle évite d'automatiser un blocage sur la seule base d'une signature pédagogique.
La validation suivra le protocole ci-dessous :
\begin{enumerate}
    \item Rejouer le scénario sans protection additionnelle et conserver l'alerte Suricata correspondante.
    \item Appliquer la règle de blocage sur la passerelle après confirmation de l'alerte.
    \item Rejouer le même scénario et constater l'échec du flux vers DVWA\@.
    \item Retirer la règle, puis vérifier le retour à l'état initial.
\end{enumerate}

Le blocage de l'adresse source protège efficacement cette démonstration, mais reste une mesure grossière : il ne corrige pas la vulnérabilité applicative, peut bloquer un utilisateur légitime et ne remplace pas les contre-mesures propres à l'application.
